% Options for packages loaded elsewhere
% Options for packages loaded elsewhere
\PassOptionsToPackage{unicode}{hyperref}
\PassOptionsToPackage{hyphens}{url}
\PassOptionsToPackage{dvipsnames,svgnames,x11names}{xcolor}
%
\documentclass[
  letterpaper,
  DIV=11,
  numbers=noendperiod]{scrartcl}
\usepackage{xcolor}
\usepackage{amsmath,amssymb}
\setcounter{secnumdepth}{-\maxdimen} % remove section numbering
\usepackage{iftex}
\ifPDFTeX
  \usepackage[T1]{fontenc}
  \usepackage[utf8]{inputenc}
  \usepackage{textcomp} % provide euro and other symbols
\else % if luatex or xetex
  \usepackage{unicode-math} % this also loads fontspec
  \defaultfontfeatures{Scale=MatchLowercase}
  \defaultfontfeatures[\rmfamily]{Ligatures=TeX,Scale=1}
\fi
\usepackage{lmodern}
\ifPDFTeX\else
  % xetex/luatex font selection
\fi
% Use upquote if available, for straight quotes in verbatim environments
\IfFileExists{upquote.sty}{\usepackage{upquote}}{}
\IfFileExists{microtype.sty}{% use microtype if available
  \usepackage[]{microtype}
  \UseMicrotypeSet[protrusion]{basicmath} % disable protrusion for tt fonts
}{}
\makeatletter
\@ifundefined{KOMAClassName}{% if non-KOMA class
  \IfFileExists{parskip.sty}{%
    \usepackage{parskip}
  }{% else
    \setlength{\parindent}{0pt}
    \setlength{\parskip}{6pt plus 2pt minus 1pt}}
}{% if KOMA class
  \KOMAoptions{parskip=half}}
\makeatother
% Make \paragraph and \subparagraph free-standing
\makeatletter
\ifx\paragraph\undefined\else
  \let\oldparagraph\paragraph
  \renewcommand{\paragraph}{
    \@ifstar
      \xxxParagraphStar
      \xxxParagraphNoStar
  }
  \newcommand{\xxxParagraphStar}[1]{\oldparagraph*{#1}\mbox{}}
  \newcommand{\xxxParagraphNoStar}[1]{\oldparagraph{#1}\mbox{}}
\fi
\ifx\subparagraph\undefined\else
  \let\oldsubparagraph\subparagraph
  \renewcommand{\subparagraph}{
    \@ifstar
      \xxxSubParagraphStar
      \xxxSubParagraphNoStar
  }
  \newcommand{\xxxSubParagraphStar}[1]{\oldsubparagraph*{#1}\mbox{}}
  \newcommand{\xxxSubParagraphNoStar}[1]{\oldsubparagraph{#1}\mbox{}}
\fi
\makeatother


\usepackage{longtable,booktabs,array}
\newcounter{none} % for unnumbered tables
\usepackage{calc} % for calculating minipage widths
% Correct order of tables after \paragraph or \subparagraph
\usepackage{etoolbox}
\makeatletter
\patchcmd\longtable{\par}{\if@noskipsec\mbox{}\fi\par}{}{}
\makeatother
% Allow footnotes in longtable head/foot
\IfFileExists{footnotehyper.sty}{\usepackage{footnotehyper}}{\usepackage{footnote}}
\makesavenoteenv{longtable}
\usepackage{graphicx}
\makeatletter
\newsavebox\pandoc@box
\newcommand*\pandocbounded[1]{% scales image to fit in text height/width
  \sbox\pandoc@box{#1}%
  \Gscale@div\@tempa{\textheight}{\dimexpr\ht\pandoc@box+\dp\pandoc@box\relax}%
  \Gscale@div\@tempb{\linewidth}{\wd\pandoc@box}%
  \ifdim\@tempb\p@<\@tempa\p@\let\@tempa\@tempb\fi% select the smaller of both
  \ifdim\@tempa\p@<\p@\scalebox{\@tempa}{\usebox\pandoc@box}%
  \else\usebox{\pandoc@box}%
  \fi%
}
% Set default figure placement to htbp
\def\fps@figure{htbp}
\makeatother





\setlength{\emergencystretch}{3em} % prevent overfull lines

\providecommand{\tightlist}{%
  \setlength{\itemsep}{0pt}\setlength{\parskip}{0pt}}



 


\KOMAoption{captions}{tableheading}
\makeatletter
\@ifpackageloaded{caption}{}{\usepackage{caption}}
\AtBeginDocument{%
\ifdefined\contentsname
  \renewcommand*\contentsname{Table of contents}
\else
  \newcommand\contentsname{Table of contents}
\fi
\ifdefined\listfigurename
  \renewcommand*\listfigurename{List of Figures}
\else
  \newcommand\listfigurename{List of Figures}
\fi
\ifdefined\listtablename
  \renewcommand*\listtablename{List of Tables}
\else
  \newcommand\listtablename{List of Tables}
\fi
\ifdefined\figurename
  \renewcommand*\figurename{Figure}
\else
  \newcommand\figurename{Figure}
\fi
\ifdefined\tablename
  \renewcommand*\tablename{Table}
\else
  \newcommand\tablename{Table}
\fi
}
\@ifpackageloaded{float}{}{\usepackage{float}}
\floatstyle{ruled}
\@ifundefined{c@chapter}{\newfloat{codelisting}{h}{lop}}{\newfloat{codelisting}{h}{lop}[chapter]}
\floatname{codelisting}{Listing}
\newcommand*\listoflistings{\listof{codelisting}{List of Listings}}
\makeatother
\makeatletter
\makeatother
\makeatletter
\@ifpackageloaded{caption}{}{\usepackage{caption}}
\@ifpackageloaded{subcaption}{}{\usepackage{subcaption}}
\makeatother
\usepackage{bookmark}
\IfFileExists{xurl.sty}{\usepackage{xurl}}{} % add URL line breaks if available
\urlstyle{same}
\hypersetup{
  pdftitle={Feedback on Assignment: Strengths and Weaknesses},
  colorlinks=true,
  linkcolor={blue},
  filecolor={Maroon},
  citecolor={Blue},
  urlcolor={Blue},
  pdfcreator={LaTeX via pandoc}}


\title{Feedback on Assignment: Strengths and Weaknesses}
\author{}
\date{}
\begin{document}
\maketitle

\renewcommand*\contentsname{Table of contents}
{
\setcounter{tocdepth}{3}
\tableofcontents
}
\listoffigures
\listoftables

\subsection{Strengths}\label{strengths}

\subsubsection{1. Clear and relevant research
question}\label{clear-and-relevant-research-question}

\begin{itemize}
\tightlist
\item
  The topic ``Does higher education lead to higher salary?'' is
  well-motivated and clearly connected to labor economics and relevant
  literature (e.g.~Blaug, 1972).\\
\item
  The data choice (large Kaggle job-salary dataset) is appropriate for
  the question.
\end{itemize}

\subsubsection{2. Good coverage of required econometric
components}\label{good-coverage-of-required-econometric-components}

The assignment touches on all the requested elements:

\begin{itemize}
\tightlist
\item
  \textbf{Simple regression}: You estimate a bivariate model (salary on
  education), present coefficients, p‑values, \(R^2\), and a prediction
  interval.\\
\item
  \textbf{Multiple regression}: You add experience and industry,
  interpret coefficients, compare with the simple model, and correctly
  discuss omitted-variable bias.\\
\item
  \textbf{Dummy-variable model}: You construct and interpret remote‑work
  dummies (( \text{remote} = \text{YES}, \text{NO} ), hybrid as base),
  and qualitatively discuss their effect.\\
\item
  \textbf{Model comparison}: You compute predicted salaries for specific
  profiles in different models and argue for a preferred model using
  \(R^2\) and interpretation.
\end{itemize}

This shows you understand the basic toolbox: coefficients, p‑values,
\(R^2\), prediction, and dummy coding.

\subsubsection{3. Interpretation of coefficients and
significance}\label{interpretation-of-coefficients-and-significance}

\begin{itemize}
\tightlist
\item
  You consistently interpret intercepts and slopes in meaningful
  economic terms (e.g.~``predicted yearly salary for lowest education''
  and ``additional education decreases salary by 482.9'').\\
\item
  You correctly use the 5\% significance level and recognize that very
  small p‑values imply strong statistical significance.\\
\item
  In the multiple regression, you identify that experience has a large
  positive effect and industry is not significant, and you explicitly
  connect that to the p‑values.
\end{itemize}

\subsubsection{4. Awareness of model
limitations}\label{awareness-of-model-limitations}

\begin{itemize}
\tightlist
\item
  You notice the very low \(R^2\) in the simple regression and correctly
  conclude that education alone explains almost none of the variation in
  salary.\\
\item
  You question the negative coefficient on education and explicitly
  mention omitted variables such as work experience, skills, and
  industry.\\
\item
  In the multiple regression, you explicitly mention possible omitted
  variable bias in the simple model and recognize that including
  experience changes the education effect.
\end{itemize}

\subsubsection{5. Use of prediction and
intervals}\label{use-of-prediction-and-intervals}

\begin{itemize}
\tightlist
\item
  You go beyond just presenting coefficients and actually \textbf{use
  the model} for:

  \begin{itemize}
  \tightlist
  \item
    Point predictions for several education levels.
  \item
    A (approximate) 95\% prediction interval for a specific education
    level.
  \end{itemize}
\item
  This shows good understanding of how regression is used in practice,
  not merely as an estimation exercise.
\end{itemize}

\subsubsection{6. Structure roughly follows assignment
instructions}\label{structure-roughly-follows-assignment-instructions}

\begin{itemize}
\tightlist
\item
  Sections 2, 3, 4, and 5 are clearly labelled in line with the task
  (simple regression, multiple regression, dummy analysis, model
  selection).
\item
  You also link to the script/estm.R and mention where the dataset is
  stored, which is good practice for reproducibility.
\end{itemize}

\begin{center}\rule{0.5\linewidth}{0.5pt}\end{center}

\subsection{Weaknesses and Suggestions for
Improvement}\label{weaknesses-and-suggestions-for-improvement}

\subsubsection{1. Conceptual issues: sign and coding of ``education
level''}\label{conceptual-issues-sign-and-coding-of-education-level}

\textbf{Problem:}\\
- You find that higher education is associated with \emph{lower} salary,
both in the simple and multiple models.\\
- You report ``Education level = -482.9'' in simple regression and
``Education level = -463.991'' in multiple regression.

\textbf{Likely cause:}\\
- You probably have \textbf{reverse coding} (e.g.~1 = highest education,
5 = lowest) or a non‑intuitive numeric coding of education. In that
case, a \textbf{negative} coefficient is actually consistent with
\emph{higher} education increasing salary (because the number goes down
as education goes up).

\textbf{What's missing:}\\
- You never show or discuss \textbf{how education is coded} (which
categories, which numbers).\\
- You therefore interpret the sign mechanically as ``negative
correlation'', which may be wrong if the coding is reversed.

\textbf{Suggestion:}\\
- Explicitly describe the encoding (e.g.~1 = PhD, 2 = Master, \ldots, 5
= High school).\\
- Then re-interpret the sign correctly or, better: recode education so
that higher values mean higher education, so a positive coefficient
directly maps to your RQ.\\
- If you keep strange coding, show a small table mapping category → code
and interpret accordingly.

\subsubsection{2. Substantive interpretation: significance vs.~effect
size}\label{substantive-interpretation-significance-vs.-effect-size}

You sometimes conflate \emph{statistical significance} with
\emph{economic importance}:

\begin{itemize}
\tightlist
\item
  Simple regression: p‑value is extremely small, so you say ``there is a
  big significance between salary and education level'', but
  \(R^2 \approx 0.0003\); the effect is statistically clear but
  practically tiny.\\
\item
  Multiple regression: you call experience's effect ``extremely
  positive'' without putting it in context (e.g.~what is the typical
  salary or range?).
\end{itemize}

\textbf{Suggestions:}

\begin{itemize}
\tightlist
\item
  Distinguish clearly between:

  \begin{itemize}
  \tightlist
  \item
    \textbf{Statistical significance} (p‑value \textless{} 0.05).
  \item
    \textbf{Magnitude / practical significance} (is 5000 or 2700 large
    or small relative to average salary?).\\
  \end{itemize}
\item
  When \(R^2\) is close to 0, explicitly state that ``although the
  effect is statistically significant due to large \(n\), the model has
  almost no explanatory power.''
\end{itemize}

\subsubsection{3. Prediction interval calculation is only approximate
and not transparently
justified}\label{prediction-interval-calculation-is-only-approximate-and-not-transparently-justified}

For the 95\% prediction interval:

\begin{itemize}
\tightlist
\item
  You take \(1.96 \cdot \text{SER}\) and simply add/subtract from the
  prediction: {[} \hat{y} \pm 1.96 \cdot \hat\sigma ,. {]}
\item
  You do not:

  \begin{itemize}
  \tightlist
  \item
    Explain that this ignores the \((1 + h_0)\) factor from the exact
    formula.
  \item
    State that you are ignoring estimation uncertainty in the slope and
    intercept.
  \end{itemize}
\end{itemize}

This is \textbf{acceptable as a rough approximation}, but:

\textbf{Suggestions:}

\begin{itemize}
\tightlist
\item
  Explicitly acknowledge: ``We approximate the prediction interval using
  \(\hat y \pm 1.96\hat\sigma\), effectively assuming a large sample and
  ignoring the leverage term.''\\
\item
  Better: show the formula and, if using R, use
  \texttt{predict(...,\ interval\ =\ "prediction")} and report the exact
  interval.
\end{itemize}

\subsubsection{\texorpdfstring{4. Over‑reliance on \(R^2\) for model
choice}{4. Over‑reliance on R\^{}2 for model choice}}\label{overreliance-on-r2-for-model-choice}

In your conclusion you prefer model 2 largely because:

\begin{itemize}
\tightlist
\item
  It has the highest \(R^2\) (19.2\% vs.~very low values in others).
\end{itemize}

\textbf{This is incomplete:}

\begin{itemize}
\tightlist
\item
  \(R^2\) alone is not always the best criterion; adjusted \(R^2\) or
  AIC/BIC are often better for comparing models with different numbers
  of predictors.\\
\item
  The theoretical plausibility and interpretability of the model should
  be emphasized more than just \(R^2\).
\end{itemize}

\textbf{Suggestions:}

\begin{itemize}
\tightlist
\item
  Mention at least one additional model selection criterion: adjusted
  \(R^2\), AIC, BIC, out-of-sample RMSE, etc.\\
\item
  Note that adding many weak variables can increase raw \(R^2\) but harm
  generalization.\\
\item
  Argue that model 2 is preferred because:

  \begin{itemize}
  \tightlist
  \item
    It includes theoretically important variables (education,
    experience, remote work etc.).
  \item
    Coefficients make substantive sense.
  \item
    It improves adjusted \(R^2\) compared to the simple model.
  \end{itemize}
\end{itemize}

\subsubsection{5. Incomplete description of variables and data
cleaning}\label{incomplete-description-of-variables-and-data-cleaning}

Throughout the assignment, some key points about the dataset and
variables are missing or vague:

\begin{itemize}
\tightlist
\item
  You mention 250,000 observations and 10 variables, but you do not
  provide:

  \begin{itemize}
  \tightlist
  \item
    A clear \textbf{data description} table (variable names, types,
    units, coding).\\
  \item
    Any discussion of \textbf{missing values}, outliers, or
    transformations (are salaries in yearly units? Which currency?
    Log‑transform considered?).\\
  \end{itemize}
\item
  For some variables (industry, remote work), we do not see:

  \begin{itemize}
  \tightlist
  \item
    How many categories they have.
  \item
    Which one is the base category in regression (you mention hybrid as
    reference, but for industry you just say ``not significant'').
  \end{itemize}
\end{itemize}

\textbf{Suggestions:}

\begin{itemize}
\tightlist
\item
  Add a small table in Section 1:

  \begin{itemize}
  \tightlist
  \item
    Variable, description, type (numeric, factor), codings and units.\\
  \end{itemize}
\item
  Briefly describe any cleaning steps: dropped missing salaries?
  Filtered unrealistic values?\\
\item
  For categorical predictors, always state:

  \begin{itemize}
  \tightlist
  \item
    Reference category.
  \item
    How many dummies were included.
  \end{itemize}
\end{itemize}

\subsubsection{6. Dummy variable section is
underdeveloped}\label{dummy-variable-section-is-underdeveloped}

You do use remote‑work dummies but:

\begin{itemize}
\tightlist
\item
  No \textbf{regression table} is shown (coefficients, SE, p-values,
  \(R^2\)).\\
\item
  You say ``the model does not include interaction terms,'' but the task
  \emph{explicitly} allows/encourages interaction effects.
\end{itemize}

\textbf{Suggestions:}

\begin{itemize}
\tightlist
\item
  Present the remote‑work regression output in a small table: intercept,
  Remote=YES, Remote=NO, their p‑values, \(R^2\).\\
\item
  Consider at least \textbf{one interaction}, e.g.:

  \begin{itemize}
  \tightlist
  \item
    Remote × Experience (does experience pay off more in remote jobs?),
  \item
    Remote × Education (is the education premium different when working
    remotely?).\\
  \end{itemize}
\item
  Interpret any interaction: explain how the effect of education or
  experience differs across remote/hybrid/on‑site groups.
\end{itemize}

\subsubsection{7. Formatting, LaTeX, and minor technical
issues}\label{formatting-latex-and-minor-technical-issues}

There are several small technical and presentational problems:

\begin{itemize}
\tightlist
\item
  Spelling/grammar:

  \begin{itemize}
  \tightlist
  \item
    ``higer salary'', ``higer'', ``significance is big'', ``indusrty''
    etc.\\
  \item
    Careless spelling weakens the professional impression.
  \end{itemize}
\item
  LaTeX / formatting:

  \begin{itemize}
  \tightlist
  \item
    Expressions like \texttt{\$\$6.45*10\^{}-20\$\$} should be
    \texttt{\$6.45\ \textbackslash{}cdot\ 10\^{}\{-20\}\$}.\\
  \item
    In one place you write \texttt{14-5} which is unclear; probably
    \texttt{–1.45} or similar.\\
  \item
    Some formulas repeat words (``Salery'') and mix comma/point as
    decimal separators.
  \end{itemize}
\item
  Word count code:

  \begin{itemize}
  \tightlist
  \item
    You define
    \texttt{ford\_count\ \textless{}-\ count\_words\_qmd("assignment\_r.qmd")}
    and then \texttt{cat("Word\ count:",\ word\_count)}, but
    \texttt{word\_count} is never defined; this code will error.
  \end{itemize}
\end{itemize}

\textbf{Suggestions:}

\begin{itemize}
\tightlist
\item
  Run a spellchecker and grammar check on the entire text.\\
\item
  Standardize LaTeX to proper math notation and use a single decimal
  separator.\\
\item
  Fix or remove non‑working code chunks (or show the actual word count,
  not broken code).\\
\item
  Clean the Quarto structure: remove duplicated
  \texttt{knitr::knit\_exit()} chunks; ensure YAML keys are correctly
  indented.
\end{itemize}

\subsubsection{8. Lack of model
diagnostics}\label{lack-of-model-diagnostics}

The assignment does not discuss:

\begin{itemize}
\tightlist
\item
  Residual plots,
\item
  Heteroskedasticity,
\item
  Nonlinearity (e.g.~log-salary, squared experience),
\item
  Multicollinearity.
\end{itemize}

Given the large dataset, violations of classical assumptions are likely.

\textbf{Suggestions:}

\begin{itemize}
\tightlist
\item
  Briefly inspect at least one diagnostic, e.g.:

  \begin{itemize}
  \tightlist
  \item
    Plot residuals vs fitted.
  \item
    Histogram or QQ-plot of residuals.\\
  \end{itemize}
\item
  Mention whether heteroskedasticity or nonlinearity appears to be a
  concern and, if so, suggest remedies (robust SEs, log-salary,
  polynomial terms).
\end{itemize}

\begin{center}\rule{0.5\linewidth}{0.5pt}\end{center}

\subsection{Summary}\label{summary}

Overall, the assignment successfully:

\begin{itemize}
\tightlist
\item
  Answers the research question using multiple models,
\item
  Demonstrates understanding of regression coefficients, significance,
  prediction, and dummy variables,
\item
  Recognizes omitted-variable bias and the importance of including
  experience.
\end{itemize}

However, for a stronger econometrics mini‑thesis, you should:

\begin{itemize}
\tightlist
\item
  Clarify and possibly recode key variables (especially education) to
  ensure the sign and interpretation make economic sense,\\
\item
  Separate statistical significance from economic importance,\\
\item
  Use more rigorous model diagnostics and selection criteria,\\
\item
  Develop the dummy-variable and interaction analysis more fully,\\
\item
  Improve clarity of data description, coding, and formatting.
\end{itemize}

Addressing these points would make the work more robust, transparent,
and convincing both statistically and substantively.




\end{document}
